\section{3D 数据集}

\subsection{数据集的演进}

在深度学习时代之前，3D 数据集很小。Princeton Shape Benchmark 仅有 1800 个模型、180 个类别（平均每类 10 个模型）。

\begin{figure}[H]
\centering
\includegraphics[width=0.95\textwidth]{slides-images/slide-018.jpg}
\caption{3D 数据集演进：从 Princeton Shape Benchmark（1800模型）到 ShapeNet（5万模型）再到 Objaverse-XL（1000万模型）\protect\footnotemark}
\end{figure}
\footnotetext{Slides 第18页。}

关键里程碑：
\begin{itemize}
    \item \textbf{ShapeNet}（斯坦福主导）：300 万模型总量，常用子集 ShapeNet Core 有 5 万模型、55 个类别
    \item \textbf{Objaverse / Objaverse-XL}（AI2）：100 万至 1000 万模型，涵盖更多类别和更高质量纹理
    \item \textbf{CO3D}（Meta + Oxford）：真实物体的 360 度视频，约 19000 个物体
    \item \textbf{ScanNet}：真实室内场景扫描，约 1500--3000 个房间
\end{itemize}

\begin{warningbox}{3D 数据远少于 2D 数据}
即使最大的 3D 数据集（$\sim$1000万模型）也远小于 2D 图像数据集（LAION-5B 有 50 亿图像）。这个巨大的数据量差距是 3D 视觉面临的核心挑战之一，也是为什么越来越多的方法试图利用 2D 基础模型的先验知识来辅助 3D 理解。
\end{warningbox}

\subsection{3D 视觉的核心任务}

\begin{figure}[H]
\centering
\includegraphics[width=0.95\textwidth]{slides-images/slide-020.jpg}
\caption{3D 视觉任务总览：生成建模、条件生成（语言/图像）、3D 重建、形状补全、判别任务（分类/分割）、2D-3D 联合建模\protect\footnotemark}
\end{figure}
\footnotetext{Slides 第20页。}

\subsection{本章小结}

3D 数据集从千级规模发展到千万级，但仍远小于 2D 数据。利用大规模 2D 基础模型的先验知识来弥补 3D 数据不足，正成为主流研究方向。

%% ========== 3D 深度学习方法 ==========

